\begin{titlepage}
	\definecolor{forestgreen}{RGB}{34,92,56}   % 森林绿主题色
	\begin{tikzpicture}[remember picture, overlay]
		% ---- 1. 墨绿渐变铺满整个封面 ----
		\shade[top color=forestgreen, middle color=forestgreen!90, bottom color=forestgreen!60!black]
			(current page.south west) rectangle (current page.north east);

		% ---- 2. 细点阵（模拟网格，微弱纹理）----
		\foreach \x in {1,...,20} {
			\foreach \y in {1,...,29} {
				\fill[white, opacity=0.05] (\x,\y) circle[radius=0.018cm];
			}
		}

		% ---- 3. 主标题后方的柔光晕 ----
		\fill[white, opacity=0.05] ([yshift=-8.2cm]current page.north) circle[radius=7.0cm];
		\fill[white, opacity=0.05] ([yshift=-8.2cm]current page.north) circle[radius=4.6cm];

		% ---- 4. 右上：空间曲线与 Frenet 标架（曲线 = 螺旋线的侧视投影）----
		\begin{scope}[shift={([xshift=-6.2cm,yshift=-4.7cm]current page.north east)}]
			% 空间曲线 C
			\draw[white, line width=1.1pt, opacity=0.9, domain=0:6.0, samples=160, smooth]
				plot (\x, {0.9*sin(deg(1.6*\x))});
			% 点 P 处的密切圆（曲率圆，虚线）
			\draw[white, dashed, opacity=0.40, line width=0.6pt]
				(1.48,1.42) circle[radius=0.64cm];
			% 点 P 与 Frenet 标架：T（切向，实线）、N（法向，青色）、B（副法向，虚线）
			\fill[cyan!65!white, opacity=0.95] (1.2,0.846) circle[radius=0.045cm];
			\draw[white, line width=1.0pt, opacity=0.9, -{Stealth}] (1.2,0.846) -- (2.14,0.38);
			\draw[cyan!65!white, line width=1.0pt, opacity=0.95, -{Stealth}] (1.2,0.846) -- (1.58,1.61);
			\draw[white, dashed, line width=0.8pt, opacity=0.75, -{Stealth}] (1.2,0.846) -- (0.79,1.11);
			\node[anchor=west, text=white, opacity=0.9, font=\footnotesize\itshape] at (2.35,0.30) {$\mathbf{T}$};
			\node[anchor=south, text=cyan!65!white, font=\footnotesize\itshape] at (1.62,1.75) {$\mathbf{N}$};
			\node[anchor=east, text=white, opacity=0.85, font=\footnotesize\itshape] at (0.50,1.28) {$\mathbf{B}$};
		\end{scope}

		% ---- 5. 左下：球面上的曲线（单位球面 + 大圆曲线 C）----
		\begin{scope}[shift={([xshift=0.9cm,yshift=0.9cm]current page.south west)}]
			% 深度轴 y（虚线，暗示第三维）
			\draw[white, opacity=0.5, dashed, line width=0.6pt] (0,0) -- (-1.3,-1.15);
			% 坐标轴 x、z
			\draw[white, opacity=0.55, line width=0.8pt, -{Stealth}] (0,0) -- (2.7,0);
			\draw[white, opacity=0.55, line width=0.8pt, -{Stealth}] (0,0) -- (0,2.7);
			% 球面：轮廓圆 + 赤道 + 两条经线
			\draw[white, line width=1.0pt, opacity=0.85] (0,0) circle (1.9cm);
			\draw[white, dashed, opacity=0.40, line width=0.6pt] (0,0) ellipse (1.9cm and 0.75cm);
			\draw[white, dashed, opacity=0.40, line width=0.6pt, rotate=45] (0,0) ellipse (1.9cm and 0.75cm);
			\draw[white, dashed, opacity=0.40, line width=0.6pt, rotate=-45] (0,0) ellipse (1.9cm and 0.75cm);
			% 球面上的曲线 C（倾斜大圆，青色）
			\draw[cyan!65!white, line width=1.1pt, opacity=0.90, rotate=20]
				(0,0) ellipse (1.9cm and 0.6cm);
			% 曲线上一点
			\fill[cyan!65!white, opacity=0.95] (1.58,1.21) circle[radius=0.045cm];
			\node[anchor=south west, text=cyan!65!white, font=\footnotesize\itshape] at (2.05,1.35) {$C$};
			% 标注 O、x、y、z
			\node[anchor=west, text=white, opacity=0.7, font=\footnotesize\itshape] at (0.18,-0.30) {$O$};
			\node[anchor=west, text=white, opacity=0.7, font=\footnotesize\itshape] at (2.80,0.02) {$x$};
			\node[anchor=south, text=white, opacity=0.7, font=\footnotesize\itshape] at (-1.55,-1.35) {$y$};
			\node[anchor=south, text=white, opacity=0.7, font=\footnotesize\itshape] at (0.02,2.82) {$z$};
		\end{scope}

		% ---- 6. 金色细边框 ----
		\draw[gold, line width=0.7pt, opacity=0.55]
			([shift={(0.85,0.85)}]current page.south west) rectangle
			([shift={(-0.85,-0.85)}]current page.north east);

		% ---- 7. 顶部眉题（英文小字 + 金色短线）----
		\coordinate (T) at ([yshift=-3.5cm]current page.north);
		\draw[gold, line width=1.0pt] (T) ++(-3.4cm,0) -- ++(1.0cm,0);
		\draw[gold, line width=1.0pt] (T) ++(2.4cm,0) -- ++(1.0cm,0);
		\node[anchor=north, text=white, opacity=0.85, align=center] at (T) {
			{\sffamily\fontsize{13}{16}\selectfont
				\uppercase{Differential \ Geometry}}
		};

		% ---- 8. 主标题（居中，使用 \notename）----
		\node[anchor=center, align=center, text=white] at ([yshift=-8.0cm]current page.north) {%
			{\fontsize{40}{50}\sffamily\bfseries \notename}
		};

		% ---- 9. 金色双细线 + 曲率圆纹饰（同心圆 + 中心点）----
		\coordinate (C) at ([yshift=-11.4cm]current page.north);
		\draw[gold, line width=1.1pt] (C) ++(-2.2cm,0) -- ++(1.55cm,0);
		\draw[gold, line width=1.1pt] (C) ++(0.65cm,0) -- ++(1.55cm,0);
		\fill[gold, opacity=0.9] (C) circle[radius=0.05cm];
		\draw[gold, line width=0.9pt] (C) circle[radius=0.14cm];
		\draw[gold, line width=0.9pt, opacity=0.7] (C) circle[radius=0.24cm];

		% ---- 10. 副标题 ----
		\node[anchor=north, text=white, opacity=0.9, align=center]
			at ([yshift=-12.2cm]current page.north) {
			{\fontsize{16}{20}\sffamily 学\ \ 习\ \ 笔\ \ 记}
		};

		% ---- 11. 底部作者、日期信息 ----
		\coordinate (B) at ([yshift=3.4cm]current page.south);
		\draw[gold, line width=0.8pt, opacity=0.7] (B) ++(-1.1cm,0.85cm) -- ++(2.2cm,0);
		\node[anchor=south, text=white, align=center] at (B) {
			\begin{tabular}{c}
				{\fontsize{20}{24}\sffamily\bfseries \noteauthor} \\[0.4cm]
				{\large \sffamily \today}
			\end{tabular}
		};
	\end{tikzpicture}
\end{titlepage}
